\documentclass[11pt,a4paper]{article}
% Preámbulo compartido para los modelos del segundo parcial de ACN.
% Replica el formato del parcial real 2P-2S-2025.
\usepackage[utf8]{inputenc}
\usepackage[T1]{fontenc}
\usepackage[spanish]{babel}
\usepackage{amsmath,amssymb}
\usepackage[margin=2.6cm]{geometry}
\usepackage{enumitem}
\usepackage{array}

\pagestyle{plain}
\setlength{\parindent}{0pt}
\setlength{\parskip}{5pt}

% Listas de incisos: a), b), c), ...
\setlist[enumerate,1]{label=\alph*),leftmargin=2.2em,itemsep=7pt,topsep=7pt}
% Listas numeradas 1., 2., 3., ...
\newlist{puntos}{enumerate}{1}
\setlist[puntos]{label=\arabic*.,leftmargin=2.2em,itemsep=7pt,topsep=7pt}
% Listas de condiciones del enunciado
\setlist[itemize,1]{label=$\bullet$,leftmargin=2em,itemsep=5pt,topsep=6pt}

% Encabezado del parcial
\newcommand{\encabezado}[1]{%
  \begin{center}
    \textbf{\large Aplicaciones Computacionales en Negocios}\\[3pt]
    \textbf{\large Segundo semestre de 2025 --- Segundo parcial}\\[3pt]
    \textbf{\normalsize Modelo de práctica #1}
  \end{center}
  \vspace{4pt}
  \hrule height 0.8pt
  \vspace{14pt}
  Nombre:\par\vspace{12pt}
  Legajo:\par\vspace{12pt}
  DNI:\par\vspace{12pt}
  ¿Te recibís con este parcial?:
  \vspace{10pt}
  \hrule height 0.8pt
  \vspace{16pt}
}

% Título de ejercicio
\newcommand{\ejercicio}[1]{\par\medskip\noindent\textbf{Ejercicio #1.}\ }

% Atajos de notación usados por la cátedra
\newcommand{\R}{\mathbb{R}}
\newcommand{\Z}{\mathbb{Z}}
\newcommand{\dem}{\mathrm{dem}}
\newcommand{\capa}{\mathrm{cap}}
\newcommand{\costo}{\mathrm{costo}}


\begin{document}
\encabezado{5}

%======================================================================
\ejercicio{1}
Una empresa de pagos debe procesar un \emph{batch} de transacciones en la
nube. Sea $A$ el conjunto de transacciones y sea $S$ el
conjunto de servidores disponibles. Cada transacción $i \in A$ tiene un tiempo
de procesamiento $t_i \in \R_+$, y cada servidor $j \in S$ tiene un costo fijo de
activación $s_j \in \R_+$ y un costo variable $c_j \in \R_+$ por unidad de tiempo
de procesamiento. El \emph{batch} completo debe procesarse en a lo sumo $T$
unidades de tiempo. Consideremos el siguiente modelo:

\begin{align}
\min \quad & \sum_{i \in A} \sum_{j \in S} c_j\, t_i\, x_{ij}
           + \sum_{j \in S} s_j\, y_j \\
\sum_{j \in S} x_{ij} &= 1 && \forall\, i \in A \\
\sum_{i \in A} t_i\, x_{ij} &\le T\, y_j && \forall\, j \in S \\
x_{ij} &\in \{0,1\} && \forall\, i \in A,\ \forall\, j \in S \notag \\
y_j &\in \{0,1\} && \forall\, j \in S \notag
\end{align}

donde la variable binaria $x_{ij}$ indica si la transacción $i$ se procesa en el
servidor $j$, y la variable binaria $y_j$ indica si se activa el servidor $j$.

\begin{enumerate}
  \item La restricción (3) cumple simultáneamente dos funciones. ¿Cuáles? Si
  $T$ fuera arbitrariamente grande, ¿cuál sería la solución óptima? ¿Haría
  falta programación lineal entera para encontrarla?

  \item Mostrar que las desigualdades $x_{ij} \le y_j$ ($i \in A$, $j \in S$)
  son válidas. ¿Cambia su agregado el conjunto de soluciones enteras factibles?
  Exhibir una solución de la relajación lineal que las viole y explicar por qué
  agregarlas puede acelerar la resolución.

  \item Supongamos que todos los servidores son idénticos ($s_j = s$ y $c_j = c$
  para todo $j \in S$). ¿Por qué esto genera muchas soluciones óptimas
  equivalentes, y por qué eso perjudica al \emph{branch and bound}? Proponer restricciones que eliminen parte de esa simetría sin
  eliminar todas las soluciones óptimas.

  \item Supongamos ahora que tenemos un presupuesto $B$ y que queremos
  minimizar la diferencia de carga entre el servidor activo más cargado y el
  activo menos cargado (los servidores no activados no intervienen en el
  mínimo). ¿Cómo se modifica el modelo?
\end{enumerate}

\newpage
%======================================================================
\ejercicio{2}
Tenemos que diseñar el \emph{fixture} de una liga de vóley con un conjunto
$E = \{1,\dots,n\}$ de equipos ($n$ es par) que juegan todos contra todos a una
sola rueda, en un conjunto $F = \{1,\dots,n-1\}$ de fechas. Cada equipo juega
exactamente un partido por fecha. Para cada par de equipos $i,j \in E$,
$i \ne j$, y cada fecha $f \in F$, la variable binaria $x_{ijf}$ toma el valor 1
si el equipo $i$ juega de local contra el equipo $j$ en la fecha $f$.

\begin{puntos}
  \item Escribir las restricciones básicas del modelo. Observar la diferencia
  con el modelo de \emph{double round-robin}: ¿qué restricción cambia y por qué?
  Agregar una restricción para que cada equipo juegue la misma cantidad de
  partidos de local que de visitante, con una diferencia de a lo sumo un
  partido.

  \item Agregar restricciones para que (i) ningún equipo juegue tres partidos
  consecutivos con la misma localía, y (ii) ningún equipo juegue de visitante
  simultáneamente en la primera y en la última fecha.

  \item En cada fecha se televisa exactamente un partido, que debe ser uno de los
  partidos atractivos del conjunto $\mathcal{A} \subseteq E \times E$ (no importa
  la localía). Cada equipo debe aparecer en televisión al menos una vez y a lo
  sumo tres veces a lo largo del torneo. Agregar las variables y restricciones
  necesarias para modelar estas condiciones.

  \item Para reducir la dificultad del modelo, se propone resolverlo en etapas a
  partir de un conjunto $\mathcal{P}$ de patrones de localía. Para cada fecha
  $f \in F$, sea $H(f) \subseteq \mathcal{P}$ el conjunto de patrones que juegan
  de local en la fecha $f$, y sea $z_{ip}$ una variable binaria que indica si el
  equipo $i$ usa el patrón $p$. Escribir las restricciones que asignan un patrón
  a cada equipo y que ligan los patrones con las variables $x_{ijf}$. Explicar
  por qué, en un torneo a una sola rueda, dos equipos no pueden usar el mismo
  patrón, y agregar esta condición. ¿Qué otra condición deben cumplir los
  patrones elegidos en cada fecha?

  \item Supongamos que, fijada una asignación de patrones a equipos obtenida en
  la etapa de selección, el modelo del \emph{fixture} resulta infactible, y que
  identificamos un subconjunto minimal $Q$ de pares (equipo, patrón) que no puede
  formar parte de ninguna solución factible. Escribir el corte que se agrega al
  modelo de selección de patrones. ¿Qué similitud tiene con una desigualdad de
  cubrimiento? ¿Por qué el proceso termina?
\end{puntos}

\newpage
%======================================================================
\ejercicio{3}
Una fábrica de alfajores planifica su producción mensual para los próximos 18
meses, $T = \{1,\dots,18\}$. La producción tiene tres niveles: primero se
fabrica \emph{relleno} (en kilos), luego se fabrican \emph{alfajores} sueltos, y
finalmente los alfajores se envasan en dos presentaciones: la \emph{caja} de 6
unidades y el \emph{pack} de 12 unidades. Cada alfajor consume $0.04$ kilos de
relleno. Llamamos $K = \{\mathrm{rel}, \mathrm{alf}, \mathrm{caja},
\mathrm{pack}\}$ al conjunto de elementos, y para cada $k \in K$ y cada
$t \in T$ conocemos la demanda $d_{kt} \in \R_+$ (se puede vender relleno a otras
fábricas y alfajores sueltos a kioscos), el costo unitario de producción
$c_{k} \in \R_+$, el costo unitario de almacenamiento $h_k \in \R_+$ y el stock
inicial $s_{k0}$. La planificación debe respetar las siguientes condiciones:

\begin{itemize}
  \item El relleno se fabrica en lotes exactos de 200 kilos.
  \item Si en un mes se fabrican alfajores, se deben fabricar al menos 5000
  unidades.
  \item Entre dos meses con producción de relleno deben pasar al menos dos meses
  sin producción de relleno.
  \item Cada elemento $k$ tiene un stock de seguridad $m_k$ que se debe mantener
  al final de cada mes.
  \item Cada elemento $k$ se fabrica en una máquina $\mu(k) \in U$ (donde $U$ es
  el conjunto de máquinas), que tiene un pool de $H_{ut}$ horas disponibles en
  el mes $t$; cada unidad del elemento $k$ consume $\tau_k$ horas de su máquina.
  Varios elementos pueden compartir una misma máquina.
\end{itemize}

\begin{puntos}
  \item Formular un modelo de programación lineal entera mixta para determinar
  cuánto fabricar de cada elemento en cada mes, minimizando los costos totales
  de producción y almacenamiento. Prestar especial atención a la restricción de
  balance de stock: la producción de un elemento consume stock de los elementos
  del nivel anterior.

  \item El relleno es perecedero y por lo tanto no puede quedar stock de relleno
  al final de ningún mes (es decir, todo el relleno fabricado en un mes se usa
  o se vende en ese mismo mes, y suponemos $m_{\mathrm{rel}} = 0$). ¿Cómo se
  agrega esta condición? Si además no
  hubiera demanda externa de relleno, ¿qué condición deberían cumplir las
  cantidades de alfajores fabricadas para que el problema sea factible? ¿Cómo
  impacta esto en el modelo?

  \item Los alfajores deben pasar un mes en cuarentena por control de calidad
  antes de poder envasarse o venderse: los alfajores fabricados en el mes $t$
  recién están disponibles en el mes $t+1$. ¿Cómo se modifica el modelo? ¿Qué
  condición sobre los stocks iniciales es necesaria para que el problema sea
  factible en el primer mes?

  \item La gerencia quiere que la carga de las máquinas sea lo más pareja
  posible a lo largo del año. Modificar la función objetivo para minimizar el
  máximo, entre todas las máquinas y todos los meses, del porcentaje de horas
  utilizadas. ¿Qué inconveniente tiene optimizar sólo este objetivo, y cómo se
  puede combinar con el objetivo de costos original?
\end{puntos}

\end{document}
